\subsection{UNIFORMITY OF COMPACT SPACES}

DEFINITION 1. A uniformity on a topological space X is said to be compatible with the topology of X if the latter coincides with the topology induced by the uniformity.

A topological space is said to be uniformizable, and its topology is said to be uniformizable, if there exists a uniformity on the space which is compatible with its topology.

There are topological spaces which are not uniformizable, for example any space which does not satisfy axiom \((\mathrm{O}_{\mathrm{III}})\) (§ 1, no. 2, Corollary 3 of Proposition 2); hence the question arises of determining under what conditions a topological space is uniformizable.

We shall not give a complete answer to this question until Chapter IX, § 1. In this section we shall examine only one important particular case, that in which X is compact. We have then the following theorem:

THEOREM 1. On a compact space X there is exactly one uniformity compatible with the topology of X; the entourages of this uniformity are all the neighbourhoods of the diagonal \(\Delta\) in \(\mathbf{X} \times \mathbf{X}\) . Furthermore, X endowed with this uniformity is a complete uniform space.

The last part of the theorem is straightforward; for every Cauchy filter on X has a cluster point {[}axiom (C){]} and is therefore convergent (§ 3, no. 2, Proposition 5, Corollary 2).

Let us show next that if there is a uniformity on X compatible with the topology of X, then the set U of entourages of this uniformity is the set of all neighbourhoods of \(\Delta\) . We know already that every entourage is a neighbourhood of \(\Delta\) (§ 1, no. 2, Proposition 2), hence we have to show that conversely every neighbourhood of \(\Delta\) belongs to U. Suppose that there is a neighbourhood U of \(\Delta\) which is not in U; then the sets \(V \cap CU\) , as V runs through U, form a base of a filter G on the compact space \(X \times X\) ; consequently G has a cluster point \((a, b)\) not belonging to \(\Delta\) . Since U is a filter coarser than G, \((a, b)\) is also a cluster point of U. But the uniformity defined by U is Hausdorff by hypothesis; hence the intersection of the closures of the sets of U is \(\Delta\) (§ 1, no. 2, Corollary 2 to Proposition 2 and Proposition 3); thus we arrive at a contradiction.

It remains therefore to show that the set \(\mathfrak{B}\) of neighbourhoods of \(\Delta\) in \(\mathbf{X} \times \mathbf{X}\) is the set of entourages of a uniformity compatible with the topology of \(\mathbf{X}\) . For this it is enough to show that \(\mathfrak{B}\) is the set of entourages of a Hausdorff uniformity on \(\mathbf{X}\) ; for then the topology induced by this uniformity will be coarser than the topology of \(\mathbf{X}\) (Chapter I, § 2, no. 2, Proposition 3) and must therefore coincide with the latter (Chapter I, § 9, no. 4, Theorem 2, Corollary 3).

\(\mathfrak{B}\) clearly satisfies axioms \((\mathbf{F}_{\mathrm{I}})\) and \((\mathbf{F}_{\mathrm{II}})\) . Let us show that axioms \((\mathbf{U}_{\mathrm{II}})\) and \((\mathbf{U}_{\mathrm{III}})\) are also satisfied and that \(\Delta\) is the intersection of the sets of \(\mathfrak{B}\) . Take the last point first: every set consisting of a single point \((x, y) \in \mathbf{X} \times \mathbf{X}\) is closed, since \(\mathbf{X}\) is Hausdorff; hence if \(x \neq y\) , the complement of \((x, y)\) in \(\mathbf{X} \times \mathbf{X}\) is a neighbourhood of \(\Delta\) . Since the symmetry \((x, y) \to (y, x)\) is a homeomorphism of \(\mathbf{X} \times \mathbf{X}\) onto itself, it follows that \(V \in B\) implies \(\overline{V} \in B\) , whence \((U_{II})\) . Finally, suppose that \((U_{III})\) is not satisfied; then there is a set \(V \in B\) such that for all \(W \in B\) the set \(\overline{W} \cap C V\) is not empty; the sets \(\overline{W} \cap C V\) , as W runs through B, therefore form a filter base on \(X \times X\) , and this filter base has a cluster point \((x, y)\) not in \(\Delta\) . Now, since X is regular (Chapter I, §9, no. 2, Corollary to Proposition 1) there are disjoint open neighbourhoods \(U_{1}\) of x and \(U_{2}\) of y, and there are closed neighbourhoods \(V_{1} \subset U_{1}\) , \(V_{2} \subset U_{2}\) of x and y respectively. Let \(U_{3} = C(V_{1} \cup V_{2})\) , and consider the neighbourhood

\[
\mathbf {W} = \bigcup_ {i = 1, 2, 3} (\mathrm{U} _ {i} \times \mathrm{U} _ {i}) \text {of} \Delta \text {in} \mathbf {X} \times \mathbf {X}.
\]

It follows immediately from these definitions that if \((u, v) \in W\) and \(u \in V_1\) (resp. \(u \in U_1\) ), then we must have \(v \in U_1\) (resp. \(v \in U_1 \cup U_3 = \complement V_2\) ); hence the neighbourhood \(V_1 \times V_2\) of \((x, y)\) in \(X \times X\) does not meet \(W\) , and we have a contradiction. This completes the proof.

Remark 1. For every finite open covering \(\Re = (\mathrm{U}_i)_{1\leqslant i\leqslant n}\) of \(\mathbf{X}\) , the set

\[
\mathrm{V} _ {\Re} = \bigcup_ {i = 1} ^ {n} \left(\mathrm{U} _ {i} \times \mathrm{U} _ {i}\right)
\]

is a neighbourhood of \(\Delta\) in \(\mathbf{X} \times \mathbf{X}\) , and these sets \(\mathrm{V}_{\Re}\) form a fundamental system of neighbourhoods of \(\Delta\) (and therefore a fundamental system of entourages of the unique uniformity on \(\mathbf{X}\) ). Let \(\mathbf{W}\) be any neighbourhood of \(\Delta\) in \(\mathbf{X} \times \mathbf{X}\) ; then for each \(x \in \mathbf{X}\) there is an open neighbourhood \(\mathbf{U}_x\) of \(x\) in \(\mathbf{X}\) such that \(\mathbf{U}_x \times \mathbf{U}_x \subset \mathbf{W}\) . Since the \(\mathbf{U}_x\) ( \(x \in \mathbf{X}\) ) form an open covering of \(\mathbf{X}\) , there exist a finite number of points \(x_i\) ( \(1 \leqslant i \leqslant n\) ) such that the \(\mathbf{U}_{x_i}\) ( \(1 \leqslant i \leqslant n\) ) form a covering \(\Re\) of \(\mathbf{X}\) . We have then \(\mathrm{V}_{\Re} \subset \mathrm{W}\) , which proves the assertion.

For this reason the unique uniformity on X is often called the uniformity of finite open coverings (cf.~Chapter IX, § 4, Exercise 17).

COROLLARY I. Every subspace of a compact space is uniformizable.

COROLLARY 2. Every locally compact space is uniformizable.

For by Alexandroff's theorem (Chapter I, § 9, no. 8, Theorem 4) a locally compact space is homeomorphic to a subspace of a compact space.

Remark 2. It can happen that there are several distinct uniformities compatible with the topology of a locally compact space.

For example, we have seen that on an infinite discrete space there is more than one distinct uniformity compatible with the discrete topology (§ 2, no. 2, Remark).

Nevertheless it is not the case that the uniqueness of the uniformity compatible with the topology of a uniformizable space is a property which characterizes compact spaces; there exist locally compact spaces which are not compact but whose uniformity is unique (Exercise 4).

THEOREM 2. Every continuous mapping \(f\) of a compact space \(X\) into a uniform space \(X'\) is uniformly continuous.

Let \(g = f \times f\) , then \(g\) is continuous on \(X \times X\) (Chapter I, § 4, no. 1, Proposition 1, Corollary 1); hence for each open entourage \(V'\) of \(X'\) , \(\overline{g}(V')\) is an open subset of \(X \times X\) , which evidently contains the diagonal. The theorem thus follows from Theorem 1, since the open entourages of \(X'\) form a fundamental system of entourages (§ 1, no. 2, Proposition 2, Corollary 2).

Under the hypotheses of Theorem 2, the restriction of \(f\) to any subspace \(A\) of \(X\) is uniformly continuous; hence (§ 3, no. 6, Theorem 2):

COROLLARY. Let A be a dense subspace of a compact space X, and let f be a mapping of A into a complete Hausdorff uniform space X'. Then f can be extended by continuity to the whole of X if and only if f is uniformly continuous.

